% GENERATED FILE. Edit canonical lessons, then run npm run generate:books.
% canonical-source-hash: fnv1a64:5712326c73ecb06b
% canonical-lessons: UR-C497-takhta, UR-C497-burada, UR-C497-randa, UR-C497-thonkna, UR-C497-retna

\chapter{The Carpenter's Bench}
\label{ch:ur-the-carpenter-s-bench}
\hlchaptermodality{497}

\begin{chapteropening}
\textbf{By the end of this chapter:} \emph{I can say how a carpenter cuts and fixes wood.}

Complete the last lesson of chapter 497: I can say how a carpenter cuts and fixes wood.
\end{chapteropening}

\section[\texorpdfstring{\ur{تختہ} (takhtā)}{تختہ (takhtā)}]{\ur{تختہ} (takhtā) — a plank; a board}

\label{lesson:UR-C497-takhta}



\begin{quote}
Before the new one: say the Urdu for to bore (a hole), then the Urdu for sandpaper.
\end{quote}

\subsection*{You'll want to know: \ur{تختہ}}
\textbf{\ur{تختہ}} — \emph{takhtā} — \textquotedblleft{}a plank; a board\textquotedblright{}.

\begin{quote}
\ur{بڑھئی} \ur{نے} \ur{ایک} \ur{لمبا} \ur{تختہ} \ur{خریدا۔}
\end{quote}

\begin{quote}
\emph{baṛhaī ne ek lambā takhtā kharīdā.}
\end{quote}

\begin{quote}
The carpenter bought a long plank.
\end{quote}

Masculine.

\begin{cousinweb}
\textbf{\ur{تختہ}} (\emph{takhtā}) is usually traced to Persian \emph{takhta}, \textquotedblleft{}a board, a plank\textquotedblright{}.
\end{cousinweb}

\begin{grammarlens}[title={What you've built}]
One more word for describing what you lived and felt.
\end{grammarlens}

\subsection*{Your turn}
\begin{itemize}
\raggedright
  \item \emph{Say it:} \emph{takhtā}
  \item \emph{Say it:} \emph{takhtā}, once more
  \item \emph{Say it:} \emph{baṛhaī ne ek lambā takhtā kharīdā.} — The carpenter bought a long plank.
  \item \emph{Recall:} say the Urdu for to bore (a hole), then the Urdu for sandpaper, then say \emph{takhtā} again
\end{itemize}

\begin{tcolorbox}[breakable,colback=teal!4,colframe=teal!35!black,title={Before you move on}]
What does \ur{تختہ} mean? (\textquotedblleft{}A plank; a board\textquotedblright{}.) Say it once more.
\end{tcolorbox}
\section[\texorpdfstring{\ur{برادہ} (burādā)}{برادہ (burādā)}]{\ur{برادہ} (burādā) — sawdust}

\label{lesson:UR-C497-burada}



\begin{quote}
Before the new one: say the Urdu for sandpaper, then the Urdu for a plank.
\end{quote}

\subsection*{You'll want to know: \ur{برادہ}}
\textbf{\ur{برادہ}} — \emph{burādā} — \textquotedblleft{}sawdust\textquotedblright{}.

\begin{quote}
\ur{بڑھئی} \ur{کی} \ur{دکان} \ur{میں} \ur{ہر} \ur{طرف} \ur{برادہ} \ur{تھا۔}
\end{quote}

\begin{quote}
\emph{baṛhaī kī dukān meṅ har taraf burādā thā.}
\end{quote}

\begin{quote}
There was sawdust all over the carpenter's shop.
\end{quote}

Masculine.

\begin{cousinweb}
\textbf{\ur{برادہ}} (\emph{burādā}) is usually traced to Persian \emph{burāda}, \textquotedblleft{}filings, sawdust\textquotedblright{}: what a saw or a file leaves behind.
\end{cousinweb}

\begin{grammarlens}[title={What you've built}]
One more word for describing what you lived and felt.
\end{grammarlens}

\subsection*{Your turn}
\begin{itemize}
\raggedright
  \item \emph{Say it:} \emph{burādā}
  \item \emph{Say it:} \emph{burādā}, once more
  \item \emph{Say it:} \emph{baṛhaī kī dukān meṅ har taraf burādā thā.} — There was sawdust all over the carpenter's shop.
  \item \emph{Recall:} say the Urdu for sandpaper, then the Urdu for a plank, then say \emph{burādā} again
\end{itemize}

\begin{tcolorbox}[breakable,colback=teal!4,colframe=teal!35!black,title={Before you move on}]
What does \ur{برادہ} mean? (\textquotedblleft{}Sawdust\textquotedblright{}.) Say it once more.
\end{tcolorbox}
\section[\texorpdfstring{\ur{رندا} (randā)}{رندا (randā)}]{\ur{رندا} (randā) — a carpenter's plane}

\label{lesson:UR-C497-randa}



\begin{quote}
Before the new one: say the Urdu for a plank, then the Urdu for sawdust.
\end{quote}

\subsection*{You'll want to know: \ur{رندا}}
\textbf{\ur{رندا}} — \emph{randā} — \textquotedblleft{}a carpenter's plane\textquotedblright{}.

\begin{quote}
\ur{اس} \ur{نے} \ur{تختے} \ur{پر} \ur{رندا} \ur{چلایا۔}
\end{quote}

\begin{quote}
\emph{us ne takhte par randā chalāyā.}
\end{quote}

\begin{quote}
He ran a plane over the plank.
\end{quote}

Masculine.

\begin{cousinweb}
\textbf{\ur{رندا}} (\emph{randā}) is usually traced to Persian \emph{randa}, \textquotedblleft{}a plane\textquotedblright{}, the carpenter's tool.
\end{cousinweb}

\begin{grammarlens}[title={What you've built}]
One more word for describing what you lived and felt.
\end{grammarlens}

\subsection*{Your turn}
\begin{itemize}
\raggedright
  \item \emph{Say it:} \emph{randā}
  \item \emph{Say it:} \emph{randā}, once more
  \item \emph{Say it:} \emph{us ne takhte par randā chalāyā.} — He ran a plane over the plank.
  \item \emph{Recall:} say the Urdu for a plank, then the Urdu for sawdust, then say \emph{randā} again
\end{itemize}

\begin{tcolorbox}[breakable,colback=teal!4,colframe=teal!35!black,title={Before you move on}]
What does \ur{رندا} mean? (\textquotedblleft{}A carpenter's plane\textquotedblright{}.) Say it once more.
\end{tcolorbox}
\section[\texorpdfstring{\ur{ٹھونکنا} (ṭhoṅknā)}{ٹھونکنا (ṭhoṅknā)}]{\ur{ٹھونکنا} (ṭhoṅknā) — to knock in; to hammer (a nail)}

\label{lesson:UR-C497-thonkna}



\begin{quote}
Before the new one: say the Urdu for sawdust, then the Urdu for a carpenter's plane.
\end{quote}

\subsection*{You'll want to know: \ur{ٹھونکنا}}
\textbf{\ur{ٹھونکنا}} — \emph{ṭhoṅknā} — \textquotedblleft{}to knock in; to hammer (a nail)\textquotedblright{}.

\begin{quote}
\ur{اس} \ur{نے} \ur{ہتھوڑے} \ur{سے} \ur{دیوار} \ur{میں} \ur{کیل} \ur{ٹھونکی۔}
\end{quote}

\begin{quote}
\emph{us ne hathauṛe se dīvār meṅ kīl ṭhoṅkī.}
\end{quote}

\begin{quote}
He hammered a nail into the wall.
\end{quote}

Past, with \ur{نے} (he or she …): \ur{اس} \ur{نے} \ur{تختے} \ur{میں} \ur{کیل} \ur{ٹھونکی}, \emph{us ne takhte meṅ kīl ṭhoṅkī}.

\begin{cousinweb}
\textbf{\ur{ٹھونکنا}} (\emph{ṭhoṅknā}) is inherited: CDIAL lists the Hindi cousin \emph{ṭhoknā}, \textquotedblleft{}to knock, to make firm\textquotedblright{}, under \emph{*ṭhōkk}, \textquotedblleft{}knock\textquotedblright{}, a form it reconstructs (CDIAL 5513).
\end{cousinweb}

\begin{grammarlens}[title={What you've built}]
One more word for describing what you lived and felt.
\end{grammarlens}

\subsection*{Your turn}
\begin{itemize}
\raggedright
  \item \emph{Say it:} \emph{ṭhoṅknā}
  \item \emph{Say it:} \emph{ṭhoṅknā}, once more
  \item \emph{Say it:} \emph{us ne hathauṛe se dīvār meṅ kīl ṭhoṅkī.} — He hammered a nail into the wall.
  \item \emph{Recall:} say the Urdu for sawdust, then the Urdu for a carpenter's plane, then say \emph{ṭhoṅknā} again
\end{itemize}

\begin{tcolorbox}[breakable,colback=teal!4,colframe=teal!35!black,title={Before you move on}]
What does \ur{ٹھونکنا} mean? (\textquotedblleft{}To knock in; to hammer (a nail)\textquotedblright{}.) Say it once more.
\end{tcolorbox}
\section[\texorpdfstring{\ur{ریتنا} (retnā)}{ریتنا (retnā)}]{\ur{ریتنا} (retnā) — to file (metal or wood)}

\label{lesson:UR-C497-retna}



\begin{quote}
Before the new one: say the Urdu for a carpenter's plane, then the Urdu for to knock in.
\end{quote}

\subsection*{You'll want to know: \ur{ریتنا}}
\textbf{\ur{ریتنا}} — \emph{retnā} — \textquotedblleft{}to file (metal or wood)\textquotedblright{}.

\begin{quote}
\ur{مستری} \ur{نے} \ur{چابی} \ur{کے} \ur{دانت} \ur{ریتے۔}
\end{quote}

\begin{quote}
\emph{mistrī ne chābī ke dānt rete.}
\end{quote}

\begin{quote}
The mechanic filed down the teeth of the key.
\end{quote}

Past, with \ur{نے} (he or she …): \ur{اس} \ur{نے} \ur{کیل} \ur{ریتی}, \emph{us ne kīl retī}.

\begin{cousinweb}
\textbf{\ur{ریتنا}} (\emph{retnā}) is inherited: CDIAL lists the Hindi cousins \emph{retī}, \textquotedblleft{}sand, a file\textquotedblright{}, and \emph{retnā}, \textquotedblleft{}to file\textquotedblright{}, under Sanskrit \emph{rētra} (CDIAL 10816).
\end{cousinweb}

\begin{grammarlens}[title={What you've built}]
One more word for describing what you lived and felt.
\end{grammarlens}

\subsection*{Your turn}
\begin{itemize}
\raggedright
  \item \emph{Say it:} \emph{retnā}
  \item \emph{Say it:} \emph{retnā}, once more
  \item \emph{Say it:} \emph{mistrī ne chābī ke dānt rete.} — The mechanic filed down the teeth of the key.
  \item \emph{Recall:} say the Urdu for a carpenter's plane, then the Urdu for to knock in, then say \emph{retnā} again
\end{itemize}

\begin{tcolorbox}[breakable,colback=teal!4,colframe=teal!35!black,title={Before you move on}]
What does \ur{ریتنا} mean? (\textquotedblleft{}To file (metal or wood)\textquotedblright{}.) Say it once more.
\end{tcolorbox}
